\section{All internal angular signs near the upper edge in growing dimension}\label{sec:near-full-all-angular-signs}

Let $d\to\infty$, $1\le r<d$, and $n/d^2\to\infty$. Put $H(w)=\sum_{s=0}^r\binom ds w^s$. For $t>0$, give the row size $S$ probabilities proportional to $\binom ds t^s$, and conditionally choose a uniform size-$S$ subset $I$. Write $J=r-S$, $h=\mathbb E J$, $\mu=r-h$, $\phi=\pi/d$, $T=nr/d$ and $D=\lfloor T\rfloor$. Define

\begin{equation*}
\begin{gathered}
g(\tau)=\mathbb E e^{-i\tau J},\quad
H_\tau=H(te^{i\tau})=H(t)e^{ir\tau}g(\tau),\\
a_s=\frac{s(d-s)}{d(d-1)},\quad
a=\mathbb E a_S,\quad a_\tau=\frac{\mathbb E[a_S e^{-i\tau J}]}{g(\tau)}.
\end{gathered}
\end{equation*}

There is an absolute $h_0>0$ such that the following holds eventually, uniformly for $0<h\le h_0$. Set

\begin{equation*}
\Psi=n\bigl(r\phi+\mathrm{Im}\log g(\phi)\bigr)
 -\frac{d-1}{2}\mathrm{Im}\log a_\phi .
\end{equation*}

At every node $\Psi=k\pi$ with integer $1\le k\le D-1$,

\begin{equation*}
(-1)^kP_n^{(r,d)}(-t^d)>0.
\end{equation*}

The logarithm of $g$ is the branch near one and that of $a_\tau$ is the branch in the right half-plane. The proof uses angular integrals throughout the near-full range, including $1\le D-k=o(d)$. It does not require asymptotic formulas for the upper coefficients.

\subsection{Tilt parameters and diagonal phases}

The successive ratios of the deficit weights decrease. The first is $q=\Pr(J=1)/\Pr(J=0)\le h/(1-h)$. Hence $\Pr(J=j)\le q^j$, $\mathbb E J^\ell=O_\ell(h)$ for every fixed $\ell\ge1$, and $h\le Cq$. Consequently $\Pr(J=0)\ge1-h$, $\Pr(J=1)\asymp h$, and $\mathrm{Var}(J)\asymp h$. Also

\begin{equation*}
a_{r-j}/a_r\le1+j,\qquad a\asymp a_r\ge 1/d.
\end{equation*}

In particular $\mathbb E[a_S^q J^\ell]=O_{q,\ell}(a^q h)$ for fixed positive integers $q,\ell$, and the same weighted estimate holds with $\mathbf1_{\{J>0\}}$ in place of $J^\ell$.

Uniformly on $[-\pi,\pi]$, $g=1+O(h)$, so $g$ stays bounded away from zero. Its fixed positive order derivatives are $O(h)$. Taylor expansion of its logarithm at zero gives $\mathrm{Im}\log g(\tau)=-h\tau+O(h|\tau|^3)$ for small $\tau$. Its centered characteristic function has second log-modulus derivative $-\mathrm{Var}(J)+O(h|\tau|)$ near zero. For far phases, the pair of atoms zero and one gives $|g(\tau)|\le\exp[-ch(1-\cos\tau)]$ by the independent-copy formula. Comparing both ranges with the value at $\phi$ gives

\begin{equation}
\log|H_\tau/H_\phi|\le -ch(\tau^2-\phi^2)\qquad(\phi\le|\tau|\le\pi),
\qquad
e^{-Ch\tau^2}\le |g(\tau)|\le1\quad (|\tau|\text{ small}).
\tag{S1}
\end{equation}

The bounds on the weighted moments of $J$ show that each of the first two derivatives of $a_\tau$ is $O(ah)$. Also $a_\tau-a=O(ah)$ uniformly in $\tau$, and its first derivative at zero is purely imaginary. Thus

\begin{equation}
\frac{\mathrm{Re}\,a_\tau}{a}=1+O(h\min(\tau^2,1)),\qquad
\frac{|\mathrm{Im}\,a_\tau|}{a}=O(h\min(|\tau|,1)),\qquad
\frac{|a_\tau|}{a}=1+O(h\min(\tau^2,1)).
\tag{S2}
\end{equation}

The cap range is chosen small enough for $\mathrm{Re}\,a_\tau$ to be positive uniformly.

\subsection{Global common-angle localization}

Use normalized Haar measure on $|z_i|=1$, $\prod_i z_i=-1$. With $R=H(t)\mathbb E\prod_{i\in I}z_i$, its integral after raising to the $n$th power equals $P_n^{(r,d)}(-t^d)$ by character orthogonality. Put $\overline z=d^{-1}\sum_i z_i$ and $D_z=\sum_i|z_i-\overline z|^2$.

At size $s$ put $m=\min(s,d-s)$. Generate a uniform subset by making $m$ disjoint random pairs from a uniform permutation, independently choosing one element from each pair, and including the next $s-m$ elements. Given the permutation, the absolute response is bounded by $e^{-X/8}$, where $X$ is the sum of squared chords of the pairs. Each chord is at most four and $\mathbb EX=2mD_z/(d-1)$. For two disjoint pairs the expected product of their squared chords is at most $d(d-1)/[(d-2)(d-3)]$ times the squared mean chord of one pair, by allowing intersections in the numerator. Thus $\mathbb EX^2\le C(\mathbb EX+(\mathbb EX)^2)$. Cauchy--Schwarz gives $\mathbb E[X/(1+X)]\ge(\mathbb EX)^2/\mathbb E[X+X^2]$ when the latter denominator is positive. With $1-e^{-x/8}\ge c x/(1+x)$ this bounds the fixed-size modulus by $1-c\min(mD_z/d,1)$, also valid when $m=0$. On averaging, retain just the size $r$ in the modulus loss. Its mass is at least $1-h$, and $\min(r,d-r)/d\ge (d-1)a_r/d$ with $a_r\asymp a$. Therefore

\begin{equation*}
|R/H(t)|\le 1-c\min(aD_z,1).
\end{equation*}

Choose a sufficiently small fixed $\eta>0$ for the local bounds below, and restrict to $D_z\le c_1\eta^2$ with a sufficiently small fixed $c_1$. The scale to compare against is $|H_\phi|^n I_0$, with

\begin{equation*}
I_0=\frac{(2\pi n a)^{-(d-1)/2}}{\sqrt d}.
\end{equation*}

On the outside part with $aD_z>1$, the integral is at most $H(t)^n e^{-cn}$. This is negligible on the principal scale since $n\gg d\log n$ under $n/d^2\to\infty$, and the principal normalization loss relative to $H(t)^n$ is $O(nh/d^2)$ in the exponent.

On the outside part with $aD_z\le1$, keep half the chord penalty for $e^{-cna\eta^2}$. Choose a pin $j$ minimizing $\sum_i|z_i-z_j|^2\le2D_z$ by averaging, and apply the other half of the penalty to the pinned squared chords. For a fixed pin, the map to the $d-1$ relative coordinates $w_i=z_i/z_j$, $i\ne j$, is onto, since one can choose $z_j^d=-1/\prod_{i\ne j}w_i$. This map is a translate of an onto homomorphism of compact tori, so it pushes normalized Haar measure to product Haar measure. Each chord factor integrates to at most $C/\sqrt{na}$ in one coordinate. A union bound over pins thus gives outside contribution at most $H(t)^n d e^{-cna\eta^2}(C/\sqrt{na})^{d-1}$. Compared to the principal scale, its logarithm is bounded by $-cna\eta^2+O(d+\log d+nh/d^2)$, tending to minus infinity because $na\gg d$ and $nh/d^2=o(na)$.

In the retained domain choose a pin minimizing $\sum_i|z_i-z_j|^2\le2D_z$ and lift all angular differences around it. Subtract their average to write $z_i=e^{i(\tau+y_i)}$, $\sum_i y_i=0$, and $\|y\|\le\eta$. The common phase $\tau$ is an odd multiple of $\phi$, taken in $(-\pi,\pi]$. This representation is unique within the small-norm domain, since any two such angle lifts differ by the same integer multiple of $2\pi$ in every coordinate, as seen by taking differences. Each chart contains a full fixed-radius ball in the zero-sum hyperplane, since $D_z\le\sum_i|e^{iy_i}-1|^2\le\|y\|^2$. In the first $d-1$ deviation coordinates the measure is $(2\pi)^{-(d-1)}dy_1\cdots dy_{d-1}$, with $y_d=-\sum_{i<d}y_i$.

\subsection{Transverse expansion and a uniform integral estimate}

Write $Q_j=\sum_i y_i^j$ and $X=\sum_{i\in I}y_i$. At a fixed size $s$, let $m_j=(s)_j/(d)_j$. Expanding the powers with $Q_1=0$ gives

\begin{equation*}
\begin{aligned}
\mathbb E[X\mid S=s]&=0,\\
\mathbb E[X^2\mid S=s]&=a_sQ_2,\\
\mathbb E[X^3\mid S=s]&=a_s\frac{d-2s}{d-2}Q_3,\\
\mathbb E[X^4\mid S=s]&=(m_1-7m_2+12m_3-6m_4)Q_4
       +3(m_2-2m_3+m_4)Q_2^2.
\end{aligned}
\end{equation*}

For example, the sums of the quartic monomials with patterns $3+1$, $2+2$, $2+1+1$ and $1+1+1+1$, with repeated exponents unlabeled, are $-Q_4$, $(Q_2^2-Q_4)/2$, $Q_4-Q_2^2/2$, and $Q_2^2/8-Q_4/4$. The coefficient of $Q_4$ is $O(a_s)$ by writing it in terms of $m_j-m_{j+1}=m_j(d-s)/(d-j)$, while $m_2-2m_3+m_4=(s)_2(d-s)_2/(d)_4=O(a_s^2)$. Consequently the real fourth moment is bounded by $C(a Q_4+a^2Q_2^2)$, using $\mathbb E a_S^2=O(a^2)$.

Tilt the size weights by $e^{-i\tau J}/g(\tau)$, whose absolute weights are bounded by a constant times the real probabilities. Since the linear term vanishes, Taylor through the cubic and the logarithm near one give

\begin{equation}
\log(R/H_\tau)
=-\tfrac12 a_\tau Q_2-\tfrac{i}{6}b_\tau Q_3
       +O(a Q_4+a^2Q_2^2),
\tag{S3}
\end{equation}

where $b_\tau$ is the tilted mean of $a_s(d-2s)/(d-2)$. The real exponential Taylor remainder is bounded by the absolute fourth moment for every real value of $X$. Also $|R/H_\tau-1|\le C a Q_2$ by the second-order Taylor bound, so the logarithm is well defined on this chart and its quadratic error is included in (S3). The weighted moment bounds give $|b_\tau|=O(a)$ and $|\mathrm{Im}\,b_\tau|=O(ah\min(|\tau|,1))$ by differentiating and by comparison with the $J=0$ term. Using $\min(|\tau|,1)\sum_i|y_i|^3\le[\min(\tau^2,1)Q_2+Q_4]/2$, the real part in (S3) is at most

\begin{equation}
-\tfrac12 a_{\rm low}Q_2+C(aQ_4+a^2Q_2^2),\qquad
a_{\rm low}=a[1-C_0h\min(\tau^2,1)]\asymp a.
\tag{S4}
\end{equation}

Fix $C_0$ large enough, then $h_0$ small enough. The same upper bound applies to the real part of the comparison exponent $-a_\tau Q_2/2$.

Here is a uniform bound for the positive exponential corrections. Put $A=na_{\rm low}\asymp na$. Under the normalized Gaussian on the hyperplane with density proportional to $e^{-A Q_2/2}$, the vector of deviations has the law of $Y_i=X_i-\overline X$ for independent $X_i\sim N(0,A^{-1})$. The vector $Y$ is independent of $\overline X$. On $Q_2(Y)\le\eta^2$, impose also $|\overline X|\le\eta$, which has probability at least $1/2$ independently. Then $|X_i|\le2\eta$, and $Q_4(Y)\le16\sum_iX_i^4$ by the fourth-power inequality and Jensen. For any fixed $C_1$, with $\eta$ sufficiently small, the one-dimensional Gaussian expectation of $e^{C_1 A X_i^4}$ restricted to this range is $1+O(A^{-1})$, by comparison with a Gaussian with half the original precision. The product is bounded since $d/A=o(1)$. This bounds the exponential moment of $C_1 n a Q_4$ on the localized ball, for any fixed constant needed below.

For the radial correction, $U=AQ_2$ has a chi-square law with $d-1$ degrees of freedom. The exponent $C_1na^2Q_2^2$ is $O(U^2/n)$. On $U\le Ld$ with fixed large $L$ this is $O(d^2/n)=o(1)$. On $Ld<U\le A\eta^2$ it is at most $\epsilon U$ for a fixed $\epsilon<1/8$ by choosing $\eta$ small. The exponential moment at $1/4$ bounds the contribution of that range by $e^{-c d}$ for $L$ sufficiently large. Cauchy--Schwarz combines the two quartic exponential bounds, and the same proof works with doubled constants.

The Gaussian for $a_{\rm low}$ has integral $I_0\exp(O(dh\min(\tau^2,1)))$, since the free-coordinate quadratic form has matrix $I+\mathbf1\mathbf1^\top$ and determinant $d$. From (S4) and the exponential bounds, the absolute contribution of any localized cluster is at most

\begin{equation}
C I_0|H_\tau|^n \exp(Cdh\min(\tau^2,1)).
\tag{S5}
\end{equation}

The exact complex Gaussian gives its leading term with an absolute error at most

\begin{equation}
C I_0|H_\tau|^n \exp(Cdh\min(\tau^2,1))\,
\left(\sqrt{\frac d{na}}+\frac d{na}+\frac{d^2}{n}\right).
\tag{S6}
\end{equation}

To see (6), apply $|e^A-e^B|\le|A-B|e^{\max(\mathrm{Re}\,A,\mathrm{Re}\,B)}$ to the two exponents and Cauchy--Schwarz to the exponential correction. Under the real comparison Gaussian, the $L^2$ norms of $Q_3,Q_4,Q_2^2$ are $O(\sqrt d\,A^{-3/2})$, $O(d A^{-2})$, and $O(d^2 A^{-2})$. The first follows from Wick's formula $\mathbb E Q_3^2=6A^{-3}(d-3+2/d)$, since the covariance has row sums zero. The others follow from marginal eighth moments and the chi-square fourth moment. Multiplying by the coefficients in (S3) gives (6). The Gaussian can be extended to the full hyperplane with an additional exponentially small error, since the chart contains a fixed-radius ball and $na\gg d$. Its integral after multiplying by the saddle value is

\begin{equation}
\frac{H_\tau^n(2\pi n a_\tau)^{-(d-1)/2}}{\sqrt d}.
\tag{S7}
\end{equation}

The factor in parentheses in (6) is $o(1)$ since $a\ge c/d$ and $n/d^2\to\infty$.

\subsection{Signs at all internal nodes}

At a node $\Psi=k\pi$, the tilted argument expansion and (S2) give

\begin{equation}
T-k=\frac{nh}{d}\left[1+O(d^{-2}+d/n)\right].
\tag{S8}
\end{equation}

Indeed $\mathrm{Im}\log g(\phi)=-h\phi+O(h\phi^3)$ and $\arg a_\phi=O(h\phi)$. The phase at $\phi$ uses $r\phi+\mathrm{Im}\log g(\phi)$ continuously, which is also the argument of $H_\phi$ between zero and $r\phi$. Since $k\le D-1$, (8) implies $nh\ge(1-o(1))d$. It also makes $T-k=O(d)$ equivalent to bounded $nh/d^2$.

If $nh/d^2\to\infty$, (S1) and (S5) make all nonprincipal clusters negligible on the principal scale. Indeed for $|\tau|\ge3\phi$ the factor $\exp(Cdh\tau^2)$ is absorbed in half the diagonal gap, because $n\gg d$ and $\tau^2\le(9/8)(\tau^2-\phi^2)$. Summing the remaining bound over odd phase labels gives $O(e^{-cnh/d^2})$. The principal conjugate pair in (7) then gives the asserted sign by (6), since its two terms have phase $\pm k\pi$.

It remains to consider a sequence on which $nh/d^2$ stays bounded. Put $L=d/\sqrt{nh}$. For a fixed large $M$ retain all odd phases $|\tau|\le M/\sqrt{nh}$. Take $M$ large enough that these include the principal phases and the discarded phases satisfy $|\tau|\ge\sqrt2\,\phi$. At a retained phase $\tau=\ell\phi$ with odd integer $\ell$, the argument of (7) can be taken as

\begin{equation*}
n\mu\tau+O(nh|\tau|^3)+O(dh|\tau|).
\end{equation*}

Subtracting $\ell$ times the principal argument gives a difference of the same error order, since $|\ell|\phi^3\le|\tau|^3$. Uniformly over retained labels it is $O_M((nh)^{-1/2}+d\sqrt{h/n})=o(1)$ by (8) and $n/d^2\to\infty$. Their signed real parts thus have the sign $(-1)^k$, because $\ell k$ has the same parity as $k$. The local errors (6) are $o(1)$ relative to the sum of their absolute leading amplitudes. In fact $dh\tau^2\le M^2d/n=o(1)$ by the retention condition, and (S2) compares the Gaussian factors uniformly to $I_0$.

On the scale $|H_\phi|^n I_0$, that positive amplitude sum is at least $c(1+L)$. If $L$ is large, use the at least $cL$ odd phases $|\tau|\le1/\sqrt{nh}$ and the lower bound in (S1). If $L$ is bounded, one principal phase suffices. For the discarded phases, (S1) and (S5) bound their total on the same scale by

\begin{equation*}
C(1+L)e^{-c'M^2}.
\end{equation*}

This follows from $\tau^2-\phi^2\ge\tau^2/2$, absorption of $dh\tau^2$ by the diagonal gap, and the Gaussian tail sum over an equally spaced lattice of spacing $2\pi\sqrt{nh}/d$ in the variable $\sqrt{nh}\,\tau$. Choosing $M$ fixed and large makes the discarded contribution smaller than the signed retained contribution. The exponential global error is already $o(1)$ on the principal scale. This proves the sign in the bounded case. Any hypothetical sequence of failed signs has a subsequence in one of these two regimes, proving the uniform eventual assertion.

\subsection{Ordering within the near-full interval}

The phase is strictly increasing with the real log-tilt $\eta_t=\log t$ in this interval for all sufficiently large parameters. Indeed

\begin{equation*}
\frac{d}{d\eta_t}\arg H_\phi
=\mathrm{Im}\frac{\mathbb E[S e^{i\phi S}]}{\mathbb E e^{i\phi S}}
=\phi\,\mathrm{Var}(J)+O(h\phi^3).
\end{equation*}

For the last equality center $J$, use its second and fourth moments in the numerator Taylor expansion, and the centered characteristic function $1+O(h\phi^2)$ in the denominator. Also $\partial_{\eta_t}a_\tau=-\mathrm{Cov}_\tau(a_S,J)$, where the covariance is under the complex size tilt. The weighted deficit moment estimates give $\partial_\tau\partial_{\eta_t}\log a_\tau=O(h)$, by differentiating this covariance and using $|a_\tau|\asymp a$. At $\tau=0$ the log-tilt derivative is real, so $(d/d\eta_t)\arg a_\phi=O(h\phi)$. Since $\mathrm{Var}(J)\asymp h$ and $n\gg d$, the derivative of $\Psi$ is positive. At $t\to\infty$ its limit is $T\pi$. Thus the near-full nodes are ordered and each target phase strictly between the initial phase on this interval and $T\pi$ is attained once.
